% Source presentations for register row N-c-024.

\subsection{Canonical databases, containment, and recursive Datalog}
\label{sec:p09-presentation-n-c-024}

\subsubsection{The Chandra--Merlin homomorphism proof}

An equality-free conjunctive query has a tuple of distinguished variables and
an existentially quantified conjunction of relational atoms:
\[
 q(\bar x)=\exists\bar y\,
   \bigwedge_{j=1}^{m}R_j(\bar z_j).
\]
Its result on a finite database \(D\) is the set of tuples \(\bar a\) for which
some assignment sends \(\bar x\) to \(\bar a\) and makes every displayed atom
a fact of \(D\).  When \(\bar x\) is empty, the query is Boolean and its result
is truth or falsity rather than a tuple relation.

For example, the Boolean query
\[
 b()=\exists x\,y\,z\,
   (E(x,y)\land E(y,z)\land Red(z))
\]
has a canonical database with elements \(x,y,z\) and facts
\(E(x,y),E(y,z),Red(z)\), with no distinguished tuple.  The non-Boolean query
\[
 r(x)=\exists y\,(E(x,y)\land Red(y))
\]
has canonical facts \(E(x,y),Red(y)\) and distinguished one-tuple \((x)\).
Declared constants, when present, are elements that every homomorphism must
fix.

In general, the canonical database \(D_q\) makes each query variable an
element, includes exactly the relational atoms of the body as facts, and
remembers the distinguished tuple \(\bar x\).  This finite structure is not the
query syntax and not an input database; it is the typed bridge between them.
An assignment witnessing \(\bar a\in q(D)\) preserves every query atom and is
therefore exactly a homomorphism
\[
 h:D_q\longrightarrow D,
 \qquad h(\bar x)=\bar a.
\]
Conversely, such a homomorphism is the required assignment.  This proves the
evaluation clause of Theorem~\ref{thm:p09-chandra-merlin}
\cite[pp.~78--82]{ChandraMerlin1977}.

The containment direction is contravariant:
\[
 q_1\subseteq q_2
 \quad\Longleftrightarrow\quad
 D_{q_2}\longrightarrow D_{q_1}.
\]
Suppose first that \(h:D_{q_2}\to D_{q_1}\) respects the distinguished tuples.
If \(\bar a\in q_1(D)\), a witness gives
\(g:D_{q_1}\to D\).  Then \(g\circ h:D_{q_2}\to D\) witnesses
\(\bar a\in q_2(D)\).  Conversely, assume \(q_1\subseteq q_2\).  The identity
map witnesses that the distinguished tuple of \(D_{q_1}\) belongs to
\(q_1(D_{q_1})\); containment places it in \(q_2(D_{q_1})\), and that query
witness is exactly a head-respecting homomorphism
\(D_{q_2}\to D_{q_1}\).  The canonical-instance step is what forces the reverse
arrow \cite[pp.~82--85]{ChandraMerlin1977}.

The reverse arrow can be checked on the smallest path example.  Let
\[
 e()=\exists x\,y\,E(x,y),
 \qquad
 p_2()=\exists x\,y\,z\,(E(x,y)\land E(y,z)).
\]
Every database satisfying \(p_2\) satisfies \(e\), so
\(p_2\subseteq e\).  Correspondingly, an edge of \(D_e\) maps to either edge
of \(D_{p_2}\): the homomorphism goes from the canonical database of the
right-hand query \(e\) to that of the left-hand query \(p_2\).  The reverse
containment is false on the database consisting of one directed edge, and no
homomorphism sends the two-edge directed path \(D_{p_2}\) into that database.

The same theorem controls minimization.  A head-preserving retraction of
\(D_q\) is a homomorphism \(r:D_q\to C\) onto a substructure
\(C\subseteq D_q\) that fixes every element of \(C\) and preserves the
distinguished tuple.  The inclusion \(C\hookrightarrow D_q\) and \(r\) are
homomorphisms in both directions, so the query represented by \(C\) is
equivalent to \(q\).  A finite canonical database with no proper such retract
is called a \emph{core}.  This name stands for a finite structure with a
distinguished tuple, not for another query language; it differs from an
arbitrary canonical database precisely by the absence of a proper
head-preserving retract and is used to state the unique minimal form.
For a finite structure, this is equivalent to saying that every
head-preserving endomorphism is an automorphism: a nonbijective endomorphism
has a power that retracts onto its proper eventual image.

For a concrete fold, compare
\[
\begin{aligned}
 q(x)&=\exists y\,z\,(E(x,y)\land E(x,z)),\\
 q'(x)&=\exists y\,E(x,y).
\end{aligned}
\]
The map identifying \(y\) and \(z\) sends \(D_q\) to \(D_{q'}\), while the
obvious inclusion sends \(D_{q'}\) to \(D_q\).  Hence the queries are equivalent
and \(q'\) is the smaller form.  Repeating proper retractions terminates because
the canonical database is finite.  If cores \(C,C'\) are homomorphically
equivalent, choose head-preserving maps \(f:C\to C'\) and \(g:C'\to C\).
The composites \(g\circ f\) and \(f\circ g\) are endomorphisms of cores and
hence automorphisms.  It follows that \(f\) and \(g\) are inverse up to those
automorphisms and therefore that \(C\cong C'\).  Thus the core query is the
unique minimal equivalent form up to variable renaming.  The fold depends on
the equality-free language: adding \(y\ne z\) would prevent the identification
and changes the theorem's scope
\cite[pp.~81--85]{ChandraMerlin1977}.

\subsubsection{Three input measures, one evaluation relation}

Let
\[
 \operatorname{Eval}(q,D,\bar a)
 \quad\Longleftrightarrow\quad
 \bar a\in q(D).
\]
Vardi's distinctions concern which arguments of this metalinguistic relation
are input, not three different query semantics \cite[pp.~137--139]{Vardi1982}.

\begin{table}[htbp]
\centering\small
\begin{tabularx}{\textwidth}{@{}l X X@{}}
\toprule
Measure & Fixed component & Varying encoded input \\
\midrule
data complexity & \(q\) & \(D,\bar a\) \\
expression complexity & \(D\) & \(q,\bar a\) \\
combined complexity & none of \(q,D\) & \(q,D,\bar a\) \\
\bottomrule
\end{tabularx}
\caption{Vardi's three evaluation input conventions.}
\label{tab:p09-query-input-measures}
\end{table}

For conjunctive queries, each fixed query has
\(\mathsf{AC}^{0}\) data evaluation, whereas Boolean combined evaluation is
\(\mathsf{NP}\)-complete.  Containment is a third, query-to-query problem and is
\(\mathsf{NP}\)-complete.  These statements are separated in
Theorems~\ref{thm:p09-cq-data-evaluation}--\ref{thm:p09-cq-containment};
the class label is never stated without what is fixed.

The upper bounds expose why the input convention matters.  Fix a conjunctive
query with \(k\) existential variables and \(m\) body atoms.  For a database of
size \(n\), membership of an output tuple is an OR over the \(n^k\) possible
witness assignments, each feeding an AND of the \(m\) required database bits.
Because \(k\) and \(m\) are constants, these are polynomial-size,
constant-depth, unbounded-fan-in circuits.  This proves the fixed-query
\(\mathsf{AC}^{0}\) data upper bound; the wiring is uniform because the query
is fixed.  When the query varies, guessing the
homomorphism \(D_q\to D\) gives the combined \(\mathsf{NP}\) upper bound.

For hardness, encode a finite graph \(G\) as the Boolean query \(q_G\) having
one variable \(x_v\) per vertex and one \(E(x_u,x_v)\) atom per represented
edge.  Evaluate it on the fixed irreflexive three-clique \(K_3\).  A witness is
exactly a graph homomorphism \(G\to K_3\), hence exactly a proper
three-coloring.  Thus even expression complexity for Boolean conjunctive
queries is \(\mathsf{NP}\)-hard, and combined evaluation is
\(\mathsf{NP}\)-complete.

Containment membership guesses and verifies the reverse homomorphism
\(D_{q_2}\to D_{q_1}\).  The same graph construction proves hardness without
confusing containment with evaluation:
\[
 q_{K_3}\subseteq q_G
 \quad\Longleftrightarrow\quad
 D_{q_G}\longrightarrow D_{q_{K_3}}
 \quad\Longleftrightarrow\quad
 G\text{ is three-colorable}.
\]
This completes the local proof of the three CQ complexity statements
\cite[pp.~80--85]{ChandraMerlin1977}.

\subsubsection{Datalog saturation and semantic coincidence}

A Datalog rule is a function-free definite clause whose head variables occur
in its positive body.  An \emph{extensional} predicate has no program rule with
that predicate in the head; its ground facts are supplied by the varying
database \(D\).  An \emph{intensional} predicate occurs in a rule head and
is derived by the fixed program.  This split is nonredundant with logical
truth: both kinds of atom belong to the same least model, but only extensional
facts vary as stored input, while intensional facts are the output of closure.

The active domain is the finite set of constants occurring in \(P\cup D\).
Ground every rule over that domain, regard the facts of \(D\) as unit clauses,
and reuse the definite-program operator from
Definition~\ref{def:p09-definite-program}:
\[
 J_0=\varnothing,
 \qquad
 J_{i+1}=T_{P\cup D}(J_i),
 \qquad
 J_\infty=\bigcup_{i<\omega}J_i.
\]
Function-freedom makes the active ground base finite.  The safety condition on
head variables ensures that every output component is bound by a positive body
match; the output cannot vary merely because the ambient domain contains an
element unmentioned by the matched data.

For a complete calculation, let the extensional database contain
\[
 E(a,b),\qquad E(b,c),\qquad E(c,d)
\]
and let the fixed Datalog program be
\[
\begin{aligned}
 Path(X,Y)&\leftarrow E(X,Y),\\
 Path(X,Z)&\leftarrow E(X,Y),Path(Y,Z).
\end{aligned}
\]
The stages are
\[
\begin{array}{c|l}
1 & D\\
2 & D\cup\{Path(a,b),Path(b,c),Path(c,d)\}\\
3 & J_2\cup\{Path(a,c),Path(b,d)\}\\
4 & J_3\cup\{Path(a,d)\}\\
5 & J_4.
\end{array}
\]
Thus the recursive relation is transitive closure on this finite instance.
It is not expressible by one ordinary conjunctive query over arbitrary finite
graphs.  To see the obstruction, suppose a fixed query with \(m\) edge atoms
defined reachability.  Its distinguished source and target must lie in the
same connected component of its canonical database; otherwise it would accept
some pair with no connecting path.  Any homomorphism of that component into a
directed path can separate the two distinguished images by at most \(m\)
edges.  But the endpoints of a directed path of length \(m+1\) are reachable.
No fixed conjunctive query therefore defines all finite transitive closures.

The semantic coincidence follows from four short arguments
\cite[chs.~12--13]{AbiteboulHullVianu1995}.

\begin{enumerate}
\item The immediate-consequence operator is monotone: enlarging an
interpretation cannot invalidate a positive rule body.
\item There are only finitely many ground intensional facts over the active
domain, so the increasing sequence reaches a fixed point after finitely many
strict additions.
\item The fixed point is a model because every satisfied rule body has already
added its head.  Every model containing the extensional database contains the
first stage and, inductively, every later stage; hence the fixed point is the
least model.
\item Induction on proof-tree height shows that each finite proof yields
membership at some stage.  Conversely, induction on the first stage at which a
fact appears builds a finite proof tree from the rule which added it.
\end{enumerate}

Thus bottom-up saturation, finite proof trees, the least fixed point, and the
least model agree.  This proof uses positivity and the finite active domain.
It is not the candidate-dependent stable-model semantics of
Definition~\ref{def:p09-stable-model}.

The fixed-program polynomial data upper bound is now immediate.  If the largest
predicate arity and the largest number of variables in a rule are the constants
\(r\) and \(v\), a database with active-domain size \(n\) has only
\(O(n^r)\) possible ground facts and \(O(n^v)\) ground instances of each fixed
rule.  Saturation adds each fact at most once and is therefore polynomial in
the data size.

A single fixed program already gives the matching \(\mathsf P\)-hardness.  An
input database can encode a monotone Boolean circuit using facts
\(One(x)\), \(And(x,y,z)\), \(Or(x,y,z)\), and \(Output(x)\), where the first
argument of a gate fact names its output and the other two name its inputs.  Use
the fixed Datalog rules
\[
\begin{aligned}
 True(X)&\leftarrow One(X),\\
 True(X)&\leftarrow Or(X,Y,Z),True(Y),\\
 True(X)&\leftarrow Or(X,Y,Z),True(Z),\\
 True(X)&\leftarrow And(X,Y,Z),True(Y),True(Z),\\
 Accept&\leftarrow Output(X),True(X).
\end{aligned}
\]
Induction along the acyclic circuit shows that \(True(g)\) is derived exactly
for gates evaluating to one; hence \(Accept\) is derived exactly when the output
gate is true.  Monotone circuit value is \(\mathsf P\)-complete under logspace
many-one reductions, so this one fixed query proves \(\mathsf P\)-complete data
evaluation.

When the program, query, and database all vary, predicate arities and rule
variable counts are no longer fixed.  The active ground base can then be
exponential in the source input, giving the explicit saturation upper bound.
The matching reduction encodes acceptance for an alternating
polynomial-space machine, using recursive intensional predicates to propagate
the existential and universal configuration conditions; since
\(\mathsf{APSPACE}=\mathsf{EXPTIME}\), it yields
\(\mathsf{EXPTIME}\)-hardness.  General Datalog-query containment is
undecidable, unlike conjunctive-query containment.  These two long reductions
remain source-attributed rather than being replaced by the local fixed-program
argument
\cite[chs.~12--16]{AbiteboulHullVianu1995}
\cite{DantsinEiterGottlobVoronkov2001}.

\begin{table}[htbp]
\centering\scriptsize
\begin{tabularx}{\textwidth}{@{}l X l X@{}}
\toprule
Language and task & Varying input & Status & Reason or boundary \\
\midrule
fixed CQ evaluation & database and output tuple & \(\mathsf{AC}^{0}\) &
fixed existential witnesses give constant-depth circuits \\
CQ evaluation & query, database, tuple & \(\mathsf{NP}\)-complete &
guess \(D_q\to D\); three-coloring already holds with \(D=K_3\) \\
CQ containment & two queries & \(\mathsf{NP}\)-complete &
guess the contravariant canonical-database homomorphism \\
fixed Datalog evaluation & extensional database & in \(\mathsf P\), and
\(\mathsf P\)-complete for some fixed program & finite saturation; one fixed
program evaluates monotone circuits \\
Datalog evaluation & program, database, query &
\(\mathsf{EXPTIME}\)-complete & combined grounding may be exponential \\
Datalog containment & two recursive queries & undecidable &
the CQ homomorphism test does not extend through unrestricted recursion \\
\bottomrule
\end{tabularx}
\caption{Evaluation and containment boundaries for conjunctive queries and
Datalog.}
\label{tab:p09-cq-datalog-boundaries}
\end{table}

Abiteboul, Hull, and Vianu supply complete definitions and the semantic proof
spines, while marking some long arguments as a \emph{crux}.  The local
monotonicity, finite-convergence, model/fixed-point, proof-tree, fixed-data
upper-bound, and fixed-program circuit arguments have been reproduced here.
The combined \(\mathsf{EXPTIME}\)-hardness and undecidable-containment
constructions remain cited at their declared source boundary; neither is
silently promoted from a sketch to a reproduced proof.
